\section{Aspectos teóricos} \label{sec:3_aspectos_teoricos}
Se plantean y describen en este apartado los aspectos teóricos y tecnológicos que sustentan 
el desarrollo del presente TFM. Estos conceptos derivan directamente del estudio del estado del 
arte y son necesarios para comprender las decisiones de diseño, los enfoques de implementación
y los criterios de evaluación que se exponen en los capítulos posteriores.

\subsection{Imágenes e Informática} \label{subsec:3_1_imagenes_informatica}
La intersección entre la Imagen y la Informática ha generado uno de los campos de estudio 
y desarrollo tencológico más dinámicos y de mayor impacto en la era digital: el 
\textit{Procesamiento Digital de Imágenes (PDI)} \cite{pdi}. Esta disciplina se basa en la adquisición,
manipulación y análisis de información visual a través de medios computacionales, transformando
el concepto de imagen de una mera representación pasiva a un flujo de datos activo y 
susceptible de tratamiento algorítmico.

Históricamente, la imagen ha sido el medio primario para el registro y la comunicación de la 
realidad. Su valor radica en la capacidad de codificar información espacial y tonal de una 
escena o un objeto, sirviendo como puente entre la percepción humana y la interpretación 
cognitiva. La imagen, en su concepción tradicional (ya sea fotográfica, pictórica u óptica),
es una función bidimensional continua que asocia un valor (intensidad o color) 
a cada punto (x, y) del plano. Dicha función se representa mátemáticamente como:

\begin{myequation}
    \begin{equation}
        I = f(x,y) \quad \text{donde } (x,y) \in \mathbb{R}^2 \text{ e } I \in \mathbb{R}
        \label{eq:3_1_imagen_funcion}
    \end{equation}
    \myequations{Representación básica de una imagen en términos matemáticos}
    \caption{Representación básica de una imagen en términos matemáticos}
\end{myequation}

En este contexto pre-digital, la calidad de la imagen dependía exclusivamente de las 
propiedades del medio de captura (emulsión química, lentes ópticas) y de la correcta 
modulación de la luz. La luminosidad y el contraste siempre han sido factores críticos, 
ya que determinan la capacidad del observador para discernir detalles y diferenciar
elementos dentro de la imagen. Un bajo contraste o una iluminación deficiente resultan en 
una pérdida irreversible de información perceptiva y analítica.

Con la adopción de tecnologías de captura electrónica y digital (cámaras
CCD, CMOS \cite{camaras_ccd_cmos}), la imagen se somete a un proceso de digitalización
que la transforma en una matriz discreta de valores numéricos, marcando el inicio de su 
tratamiento informático. Este cambio de paradigma no solo digitaliza la imagen, sino que 
también le confiere nuevas propiedades y posibilidades de manipulación en contextos 
científicos, industriales y artísticos.

\newpage

\subsection{Representación numérica de imágenes digitales} \label{subsec:3_2_representacion_numerica_imagenes_digitales}
Una vez que la imagen es capturada por un dispositivo electrónico, el proceso de digitalización
convierte la señal óptica continua en una representación numérica discreta susceptible de 
ser almacenada, transmitida y procesada por sistemas informáticos. Este proceso consta 
fundamentalmente de dos etapas: el \textit{muestreo} y la \textit{cuantificación}.

El muestreo consiste en discretizar el dominio continuo de la imagen, dividiendo 
el plano $(x, y)$ en una rejilla regular de posiciones finitas denominadas \textit{píxeles} (del inglés 
\textit{picture elements}). Cada píxel representa una pequeña región de la escena original
y constituye la unidad mínima de información espacial en una imagen digital. 
El resultado de este proceso es una imagen discreta definida sobre un conjunto finito de 
corrdenadas enteras:

\begin{myequation}
    \begin{equation}
        I[i,j] = f(i \Delta x, j \Delta y)
        \label{eq:3_2_imagen_discreta}
    \end{equation}
    \myequations{Imagen digital como función discreta muestreada}
    \caption{Imagen digital como función discreta muestreada}
\end{myequation}

donde $(i,j)$ representan las coordenadas del píxel y $\Delta x$, $\Delta y$ las resoluciones
espaciales en cada dimensión. El número total de muestras determina la resolución espacial
de la imagen y condiciona tanto su nivel de detalle como el coste computacional asociado
a su procesamiento.

La cuantificación, por su parte, transforma los valores continuos de intensidad o color capturados 
en cada píxel en un valor discreto perteneciente a un conjunto finito de niveles. Este proceso
introduce una aproximación numérica controlada por la \textit{profundidad de bits} (\textit{bit depth})
del sistema de representación. Por ejemplo, una cuantificación de 8 bits permite representar $2^8=256$ 
niveles distintos, mientras que profundidades mayores (12, 14 o 16 bits) son habituales en imágenes
crudas procedentes de sensores modernos, proporcionando una mayor fidelidad tonal y rango dinámico.

El resultado final del proceso de digitalización la Imagen Digital, que se representa como una matriz 
bidimensional ($M \times N$) de valores discretos. Cada elemento de esta matriz corresponde a un píxel
con un valor cuantificado que refleja la intensidad luminosa o el color en esa posición específica:

\begin{myequation}
    \begin{equation}
        I[m,n] \quad \text{donde } m = 0, 1, \ldots, M-1 \text{ y } n = 0, 1, \ldots, N-1
        \label{eq:3_2_imagen_matriz}
    \end{equation}
    \myequations{Imagen digital representada como matriz de píxeles}
    \caption{Imagen digital representada como matriz de píxeles}
\end{myequation}

\newpage
La \autoref{fig:3_1_muestreo_cuantificacion} ilustra el proceso de digitalización de una imagen,
mostrando la transición desde una imagen continua $f(x,y)$, pasando por la imagen muestreada $f[m,n]$,
hasta la imagen digital final $I[m,n]$.

\begin{figure}[h]
    \centering
    \begin{tikzpicture}[
        box/.style={draw=black!70, thick, minimum size=3cm},
        title/.style={font=\small, align=center}
    ]
        % BOX: Imagen continua
        \node[box, label={[title]above:Imagen continua}] (cont) at (-3,0) {};
        
        \node (image_cont) at (-3,0) {
            \includegraphics[width=3cm,height=3cm]{sections/3/img/dog.png}
        }; 
        
        \node[font=\small] at (-3, -2) {$f(x,y)$};

        % ARROW: Muestreo
        \draw[arrow, thick] (-1,0) -- node[above, font=\small] {Muestreo} (1,0);

        % BOX: Imagen muestreada
        \node[box, fill=white, label={[title]above:Imagen muestreada}] (sampled) at (3,0) {};

        \begin{scope}[shift={(3,0)}]
            \def\step{0.5}

            \foreach \x/\y/\g in {
                -1.25/ 1.25/160, -0.75/ 1.25/150, -0.25/ 1.25/140,  0.25/ 1.25/130,  0.75/ 1.25/150, 1.25/ 1.25/160,
                -1.25/ 0.75/150, -0.75/ 0.75/110,  -0.25/ 0.75/90,   0.25/ 0.75/90,   0.75/ 0.75/110, 1.25/ 0.75/150,
                -1.25/ 0.25/140, -0.75/ 0.25/90,   -0.25/ 0.25/60,   0.25/ 0.25/60,   0.75/ 0.25/90,  1.25/ 0.25/140,
                -1.25/-0.25/140, -0.75/-0.25/90,   -0.25/-0.25/50,   0.25/-0.25/50,   0.75/-0.25/90,  1.25/-0.25/140,
                -1.25/-0.75/150, -0.75/-0.75/110,  -0.25/-0.75/90,   0.25/-0.75/90,   0.75/-0.75/110, 1.25/-0.75/150,
                -1.25/-1.25/160, -0.75/-1.25/150, -0.25/-1.25/140,  0.25/-1.25/130,  0.75/-1.25/150, 1.25/-1.25/160
            }{
                \fill[gray!\g]
                    (\x-\step/2, \y-\step/2) rectangle (\x+\step/2, \y+\step/2);
            }

            \draw[black!80, thin] (-1.5, -1.5) grid[step=0.5] (1.5, 1.5);

        \end{scope}

        \node[font=\small] at (3, -2) {$f[m,n]$};
        
        % ARROW: Cuantificación
        \draw[arrow, thick] (5,0) -- node[above, font=\small] {Cuantificación} (7,0);

        % BOX: Imagen digital
        \node[box, fill=white, label={[title]above:Imagen digital}] (digital) at (9,0) {};

        \begin{scope}[shift={(9,0)}]
            \def\step{0.5}

            \foreach \x/\y/\g in {
                -1.25/ 1.25/160, -0.75/ 1.25/150, -0.25/ 1.25/140,  0.25/ 1.25/130,  0.75/ 1.25/150, 1.25/ 1.25/160,
                -1.25/ 0.75/150, -0.75/ 0.75/110,  -0.25/ 0.75/90,   0.25/ 0.75/90,   0.75/ 0.75/110, 1.25/ 0.75/150,
                -1.25/ 0.25/140, -0.75/ 0.25/90,   -0.25/ 0.25/60,   0.25/ 0.25/60,   0.75/ 0.25/90,  1.25/ 0.25/140,
                -1.25/-0.25/140, -0.75/-0.25/90,   -0.25/-0.25/50,   0.25/-0.25/50,   0.75/-0.25/90,  1.25/-0.25/140,
                -1.25/-0.75/150, -0.75/-0.75/110,  -0.25/-0.75/90,   0.25/-0.75/90,   0.75/-0.75/110, 1.25/-0.75/150,
                -1.25/-1.25/160, -0.75/-1.25/150, -0.25/-1.25/140,  0.25/-1.25/130,  0.75/-1.25/150, 1.25/-1.25/160
            }{
                \node[font=\tiny] at (\x, \y) {\g};
            }
            
            \draw[black!70, thin] (-1.5, -1.5) grid[step=0.5] (1.5, 1.5);

        \end{scope}

        \node[font=\small] at (9, -2) {$I[m,n]$};

    \end{tikzpicture}
    \caption{Proceso de digitalización de una imagen}
    \label{fig:3_1_muestreo_cuantificacion}
\end{figure}

\subsubsection{Escala de grises y color} \label{subsubsec:3_2_escala_grises_color}

\subsubsection{Espacios de color} \label{subsubsec:3_2_espacios_color}

\subsubsection{Representación de imágenes con tensores} \label{subsubsec:3_2_representacion_tensores}

\subsection{Procesamiento de imágenes en formato crudo} \label{subsec:3_3_procesamiento_imagenes_crudo}

\subsection{Fundamentos de contraste en imágenes} \label{subsec:3_4_fundamentos_contraste_imagenes} 

\subsection{Técnicas de optimización de contraste} \label{subsec:3_5_tecnicas_optimizacion_contraste}

\subsection{Procesamiento por etapas y \textit{pipelines}} \label{subsec:3_6_procesamiento_por_etapas_pipelines}

\subsection{Procesamiento vectorizado y paralelismo de datos} \label{subsec:3_7_procesamiento_vectorizado_paralelismo_datos}    

\subsection{Computación acelerada por GPU} \label{subsec:3_8_computacion_acelerada_gpu}